% 第13章 パラメータ依存性の定量分析
\section{パラメータ依存性の定量分析}
\label{ch:ch13_params}

\paragraph{この章で導くこと（入力→出力）}
各パラメータを fiducial から動かしたときの $C_\ell$ 応答を、(i) 物理機構、
(ii) 解析スケーリング、(iii) 数値（NB6 の図と\%表）の3点セットで与える。要求仕様
（README §1.4）の必須トピック $\tau$・$z_*$・$z_{\rm re}$・$\sum m_\nu$ を扱う。

\subsection{方法論}
応答は比 $C_\ell^{\rm new}/C_\ell^{\rm fid}$ または対数微分
$\partial\ln C_\ell/\partial\ln p$ で測る。何を固定するか（$\theta_*$ 固定 vs $h$ 固定）の
宣言が本質的である。本書の自前計算は \texttt{scripts/make\_param\_study.py}、
質量ニュートリノは背景を自前実装し（\texttt{massive\_nu.py}）、摂動の自由流は CLASS と
照合する。

\subsection{$A_s,n_s$}
原始スペクトル $\mathcal P_\mathcal R=A_s(k/k_{\rm pivot})^{n_s-1}$ は線形に転送されるので、
$C_\ell\propto A_s$ で全体振幅を一様にスケールする（ただし観測量は $A_se^{-2\tau}$ の組で、
$\tau$ と縮退する）。$n_s$ は $k_{\rm pivot}=0.05~\mathrm{Mpc^{-1}}$（$\ell\sim700$ 付近に
対応）まわりでスペクトルを傾け、$n_s<1$ なら大スケール（低 $\ell$）を持ち上げ小スケール
（高 $\ell$）を下げる。対数微分は $\partial\ln C_\ell/\partial n_s=\ln(k_\ell/k_{\rm pivot})$
（$k_\ell\approx\ell/\chi_*$）でシーソー状。\textbf{観測との含意}: $A_s$ 単独は $\tau$ と
縮退し、低 $\ell$ EE（再電離バンプ）で破る。

\subsection{$\Omega_bh^2$（バリオン）}
バリオンを増やすと荷重 $R=3\rho_b/4\rho_\gamma$ が上がり、(i) 振動のゼロ点シフト
$-(1+R)\Psi$ が深くなって\textbf{圧縮側（奇数）ピークが増強}、奇偶非対称
$H_1/H_2,H_3/H_2$ が拡大（第\ref{ch:ch12_peaks}章）；(ii) 音速 $c_s=c/\sqrt{3(1+R)}$ が
下がり音地平線 $r_s$ が縮むので\textbf{ピークがわずかに高 $\ell$ 側}へ；(iii) $n_e$ 増で
拡散長 $k_D^{-1}$ が縮み\textbf{減衰尾がやや高め}に残る。\textbf{含意}: 奇偶比は $\Omega_bh^2$ の
ほぼ純粋なプローブで、BBN と独立に測れる。

\subsection{$\Omega_mh^2$（物質）}
物質を増やすと等密度 $a_{\rm eq}=\Omega_{r0}/\Omega_{m0}$ が早まる。再結合が物質優勢の
奥に入るので、(i) 放射優勢期に減衰するポテンシャル $\Phi$ の駆動（radiation driving、
第\ref{ch:ch08_acoustic}章）が弱まり\textbf{第一ピーク高が下がる}；(ii) early ISW
（等密度直後の $\Phi$ 残留減衰）も減る；(iii) $\chi_*$ と $r_s$ が変わり $\theta_*$ も動く。
\textbf{含意}: 第一ピーク高（駆動）と $\theta_*$ の組で $\Omega_mh^2$ と $\Omega_bh^2$ を
分離できる。

\subsection{$\Omega_\Lambda$・曲率（幾何縮退）}
$\Lambda$ と曲率 $\Omega_k$ は再結合期の物理を変えず、\emph{距離} $\chi_*$ だけを変える。
$100\theta_*=r_s^*/\chi_*$ を一定に保つ $(\Omega_\Lambda,\Omega_k)$ の線上では一次 $C_\ell$
（$\ell\gtrsim30$）がほぼ不変——これが\textbf{幾何学的縮退}である。破れるのは (i) late ISW
（$\Lambda$ 優勢期の $\Phi$ 減衰）が効く $\ell\lesssim30$ と、(ii) レンズ・$H_0$ 直接測定・
BAO との組合せのみ。\textbf{含意}: CMB 単独では $\Omega_\Lambda$ と $\Omega_k$ は縮退し、
外部データで破る。

\subsection{光学的厚み $\tau$ と再電離 $z_{\rm re}$（導出契約 S13.1）}
再電離後の散乱で $\ell\gtrsim30$ の全ピークが一様に $e^{-2\tau}$ で抑制される。
地平線スケールを超える低 $\ell$ は抑制されない。\texttt{cmbcore} の tanh 再電離
（$z_{\rm re}=8$, $\tau=0.069$）で（図 F6.4）:
\begin{equation}
  \frac{\Delta\mathcal D_{500}}{\mathcal D_{500}}=-12.89\%
  \quad\text{vs}\quad e^{-2\tau}-1=-12.91\%,
\end{equation}
と\textbf{$0.1\%$ 以内で一致}する（仕様 F6.4 合格）。瞬間再電離では
$\tau(z_{\rm re})$ が解析式で与えられ、$z_{\rm re}=7.7\Rightarrow\tau\approx0.054$。
$A_se^{-2\tau}$ の縮退は偏光 EE 再電離バンプ（付録E、振幅 $\propto\tau^2$）で破れる。

\subsection{再結合期 $z_*$ のずれ（導出契約）}
原始ヘリウム量・微細構造定数・エネルギー注入により $z_*$ がずれる思考実験。
\texttt{Recombination(recomb\_shift)} で Peebles 係数の温度引数を一括スケールすると
（$T_b\times1.05$、図 F6.6）、
\begin{equation}
  z_*:1082\to1032,\qquad \ell_1:220\to210,
\end{equation}
すなわち $z_*$ が下がると $r_s$ が伸び $\theta_*$ が増し $\ell_1$ が低 $\ell$ へ動く
（$\ell_1\propto1/\theta_*$）。可視度幅の変化は減衰尾にも効く。

\subsection{ニュートリノ質量 $\sum m_\nu$（導出契約 S13.2）}
非相対論化は $z_{\rm nr}\simeq1890\,(m_\nu/\mathrm{eV})$ で、$\sum m_\nu\lesssim0.6$~eV
では再結合時はほぼ相対論的。一次効果は (a) $\theta_*$・$\chi_*$ を通じた背景幾何、
(b) 後期 ISW、(c) レンズ平滑化の減少、(d) 物質パワー抑制 $\Delta P/P\approx-8f_\nu$
（自由流長）。

\texttt{cmbcore} は\textbf{質量ニュートリノの背景を厳密に実装}する
（\texttt{massive\_nu.py}）。共動運動量 $q$ の Fermi--Dirac 分布
$f_0(q)=[e^{q c/k_BT_\nu}+1]^{-1}$（$T_\nu=(4/11)^{1/3}T_{\rm CMB}$）に対し、共動エネルギー
$\varepsilon=\sqrt{(qc)^2+(amc^2)^2}$ を運動量積分してエネルギー密度を得る。無次元
$x\equiv qc/k_BT_\nu$、$y(a)\equiv amc^2/k_BT_\nu$ とおくと
\begin{equation}
  \frac{\rho_{\rm ncdm}(a)}{\rho_{\rm crit,0}}
  =A\,a^{-4}\!\int_0^\infty\! x^2\sqrt{x^2+y(a)^2}\,f_0(x)\,dx,
  \qquad
  A=\frac{N_{\rm ncdm}\,(7/8)(4/11)^{4/3}\Omega_{\gamma0}}{\int_0^\infty x^3 f_0\,dx},
  \label{eq:rho_ncdm}
\end{equation}
規格化 $A$ は相対論極限（$y\to0$、$\sqrt{x^2+y^2}\to x$）で1質量種が1無質量種
$(7/8)(4/11)^{4/3}\Omega_{\gamma0}a^{-4}$ に一致するよう定めた（$\int x^3f_0=7\pi^4/120$）。
非相対論極限（$y\gg1$）では $\sqrt{x^2+y^2}\to y\propto a$ なので積分が $\propto a$ となり
$\rho_{\rm ncdm}\propto a^{-3}$（物質的）へ移行する。今日の値は
$\Omega_{\rm ncdm}h^2\approx\sum m_\nu/93~\mathrm{eV}$ を再現し、平坦条件で
$\Omega_\Lambda$ が減るため $H(z)$ と距離が変わる（$h$ 固定で
$\sum m_\nu=0.06/0.12/0.30$~eV に対し $\Omega_\Lambda{:}\,0.686\!\to\!0.684\!\to\!0.679$、
$100\theta_*{:}\,1.034\!\to\!1.040$）。$\mathcal H$ には
$E=(\mathcal H/H_0)^2$ の項 $a^2\rho_{\rm ncdm}/\rho_{\rm crit,0}$ として入る。

\begin{numreality}
背景実装は CLASS（フル ncdm）の unlensed TT 応答 $C_\ell(\sum m_\nu)/C_\ell(0)$ と
整合する（NB6 / \texttt{scripts/make\_massive\_nu.py}）。$\theta_*$ 固定の unlensed TT は
$\sum m_\nu\lesssim0.3$~eV で数\% 以内に留まり、「CMB 単独では質量に鈍く、レンズ・
LSS との組合せが要る」ことを数値で立証する。摂動レベルの自由流（$q$ グリッド階層に
よる $P(k)$ 抑制）は TT への寄与が小さく、第2版の拡張項目とする。
\end{numreality}

\subsection{まとめ: 応答行列}
パラメータ→特徴量の応答を一表にまとめると、$\theta_*$（$\to\ell_1$）はよく決まり、
$\Omega_\Lambda$/曲率や $A_se^{-2\tau}$、$\sum m_\nu$ は縮退として残る。これが
「CMB だけで何が決まり、何が縮退か」の答えである。

\paragraph{演習}
(1) NB6 で $z_{\rm re}$ を変え $\tau$ と $\ell=500$ 抑制率の関係を確かめよ。
(2) $\partial\ln\ell_1/\partial\ln z_*$ を数値から評価し解析予測と比べよ。
(3) 幾何縮退を $100\theta_*$ 一定線上で図示せよ。
